\subsubsection{Inverse Operators}\label{subsec:inverse-operators}

Recall that we say that a mapping $A:\mathcal{D}(A)\to\mathcal{R}(A)$ is one-to-one or injective if distinct points in $\mathcal{D}(A)$ get mapped to distinct points in $\mathcal{R}(A)$; i.e., if for any $x_1,x_2\in\mathcal{D}(A)$ we have
\begin{equation}\label{eq:linear-operator-injectivity-condition}
x_1\ne x_2\quad\Rightarrow\quad Ax_1\ne Ax_2.
\end{equation}


For any such mapping we can define an inverse mapping $(\mathcal{R}(A),A^{-1})$ which maps any point $y\in\mathcal{R}(A)$ to the unique point $x\in\mathcal{D}(A)$ such that $Ax=y$. This definition implies
\begin{equation}\label{eq:inverse-operator-left-identity}
A^{-1}Ax=x
\end{equation}
for every $x\in\mathcal{D}(A)$ and
\begin{equation}\label{eq:inverse-operator-right-identity}
AA^{-1}y=y
\end{equation}
for every $y\in\mathcal{R}(A)$.

The following simple but important theorem is left to the reader (Problem~\ref{ex:inverse-linear-operator-null-space-criterion}).

\begin{Theorem}
\label{thm:inverse-linear-operator-null-space}
Let $X$ and $Y$ be Banach spaces. Let $(\mathcal{D}(A),A)$ be a linear operator from $X$ to $Y$ with range $\mathcal{R}(A)$. Then in the following hold.\begin{enumerate}
\item The inverse operator $(\mathcal{R}(A),A^{-1})$ exists if and only if $\mathcal{N}(A)=\{0\}$.
\item If the inverse operator exists, it is linear.
\end{enumerate}
\end{Theorem}
